\documentclass[10pt,a4paper]{amsart}
\usepackage[margin=19mm]{geometry}
\usepackage{amsmath,amssymb,mathtools}
\usepackage[scr=rsfs,cal=euler]{mathalfa}
\usepackage{xfrac}
\usepackage[unicode,hidelinks,bookmarks=false]{hyperref}
\newcommand{\ie}{i.e.\ }
\newcommand{\cf}{cf.\ }
\newcommand{\resp}{respectively\ }
\newcommand{\Subsection}{\subsection}
\providecommand{\nobreakdash}{\nobreak\hbox{-}}
\setlength{\emergencystretch}{2em}
\multlinegap0pt
\usepackage{bm}
\usepackage{bbm}
\usepackage{dsfont}
\usepackage{xspace}
\usepackage{cancel}
\usepackage{mathtools}
\usepackage{amsthm}
\makeatletter
\@ifundefined{thm}{\newtheorem{thm}{Theorem}}{}
\@ifundefined{defn}{\newtheorem{defn}[thm]{Definition}}{}
\makeatother
\def\Aut{{\rm Aut}}
\def\FF{{\mathcal F}}
\def\H{{\rm H}}
\def\Mass{{\rm Mass}}
\def\Q{{\mathbb Q}}
\def\Val{{\nu}}
\def\bQ{{\mathbb Q}}
\def\cF{{\mathcal F}}
\def\cS{{\mathcal S}}
\def\cT{{\mathcal T}}
\def\un{{\rm un}}
\newcommand{\D}{{\rm d}}
\newcommand{\q}{{\rm q}}
\title[Malle's Conjecture for Galois octic fields over $\mathbb Q$]{Malle's Conjecture for Galois octic fields over $\mathbb Q$}
\author{Arul Shankar, Ila Varma}
\date{Source: arXiv:2505.23690v2 (CC BY 4.0)}
\begin{document}
\maketitle

\noindent\emph{Source and licence.} Malle's Conjecture for Galois octic fields over $\mathbb Q$. Version arXiv:2505.23690v2 (CC BY 4.0), distributed under the Creative Commons Attribution 4.0 International licence, as recorded in the arXiv licence metadata for this exact version and shown on the abstract page https://arxiv.org/abs/2505.23690v2. Full rights notice: licenses/arxiv-2505.23690-CC-BY-4.0.txt.\par\smallskip

\noindent\textit{Editorial scope.} Counted unit: the Thm whose source label is \texttt{th:congQuartic} in the article (source0.tex lines 2464--2473, source label \texttt{th:congQuartic}), delivered together with the article's own complete original proof of it (source0.tex lines 2474--2550, one \texttt{proof} environment of 77 lines). No mathematics is written, repaired or re-derived: statement and proof are the article's own text line for line. Required context retained uncounted: sheight (lines 690--696); localspec (lines 753--799); eqmass (lines 1901--1905); thD4quar (lines 1908--1918); thunif (lines 2417--2422). Every definition, display and label of those lines is retained unchanged. Omitted from this package and not redistributed: 1--689, 697--752, 800--1900, 1906--1907, 1919--2416, 2423--2463, 2551--2697. Nothing inside a retained excerpt is omitted or altered: every retained body is delivered exactly as the article's lines are written, so no omission receipt of any kind accompanies this package. Citations: the retained excerpts carry no citation command at all, so no bibliography entry of the article is transcribed and no reference is redistributed. Reference adaptation: none was needed and none was made; every reference inside the delivered unit resolves inside it, so no unresolved reference or citation is produced. Printed slips kept literal: no printed word, symbol, hypothesis or number of the counted statement or of its proof was altered. Layout and class: the article's own class is \texttt{article} with packages including geometry, amsmath, amsfonts, amsthm, amssymb, latexsym, amscd, amsopn, graphics, color, enumerate, lscape, enumitem, mathtools; the wrapper is the lane's pinned \texttt{amsart} editorial wrapper at 10pt on A4, which restates the theorem-like environments the retained text needs and adds the article's own macro definitions that the retained text needs (\texttt{\textbackslash Aut}, \texttt{\textbackslash D}, \texttt{\textbackslash FF}, \texttt{\textbackslash H}, \texttt{\textbackslash Mass}, \texttt{\textbackslash Q}, \texttt{\textbackslash Val}, \texttt{\textbackslash bQ}, \texttt{\textbackslash cF}, \texttt{\textbackslash cS}, \texttt{\textbackslash cT}, \texttt{\textbackslash q}, \texttt{\textbackslash un}), with extra wrapper packages (bm, bbm, dsfont, xspace, cancel, mathtools). The delivered numbering is the unit's own. Assets: no retained span contains a figure environment, a graphics-inclusion command, a PDF-page inclusion, a table or a TikZ picture; no image file is copied into this package, the package declares no asset dependency and no image file is redistributed with it. Licence: version arXiv:2505.23690v2 of this article is distributed under the Creative Commons Attribution 4.0 International licence, as recorded in the arXiv licence metadata for this exact version and shown on the abstract page https://arxiv.org/abs/2505.23690v2. A full rights notice is delivered at licenses/arxiv-2505.23690-CC-BY-4.0.txt. Characters: every retained excerpt is pure ASCII, so the delivered engine, XeLaTeX with its default Latin Modern text font, reports no missing character.
\begin{defn}\label{sheight}
Define the $S$-{\it height} $\H=\H^S$ on such pairs $(L,K)$ of a quadratic
extension $L$ of a quadratic field $K$ to be
\begin{equation*}
\H^S(L,K):=\q_\sigma^{S}(L,K)^{1/2} \cdot \q^S_{\sigma^2}(L,K) \cdot \q^S_{\tau}(L,K) \cdot\D^{S}(L,K).
\end{equation*}
\end{defn}

\begin{defn}\label{localspec}
For each place $v$ of $\Q$, we use $\Lambda_v$ to denote a nonempty set of pairs
$(L_v,K_v)$, where $K_v$ is (an isomorphism class of) a quadratic \'etale extension of $\Q_v$
and $L_v$ is (an isomorphism class of) a quadratic \'etale extension of $K_v$. We refer to such $\Lambda_v$ as a set of {\em local specifications at $v$}. 

Let $p$ be a prime. 
\begin{itemize}
    \item A pair $(L_p,K_p)$ is said to be {\it totally ramified} if both the
extensions $K_p/\Q_p$ and $L_p/K_p$ are ramified.
    \item For a prime $p$, let $\Lambda_p^{\rm ntr}$ denote the set of all pairs $(L_p,K_p)$ that are not totally ramified.
    \item A collection of local specifications
$\Lambda=(\Lambda_p)_{p \in \nu(\bQ)}$ is said to be {\it acceptable} if for all
large enough $p$, the set $\Lambda_p$ contains every pair
$(L_p,K_p)$ which is not totally ramified, i.e., $\Lambda_p^{\rm ntr} \subseteq \Lambda_p$ for large enough $p$. 
    \item  A collection $\Lambda = (\Lambda_p)_{p \in \nu(\bQ)}$ is
said to have {\it finite total ramification} if for all large enough
$p$ the set $\Lambda_p$ has no pair $(L_p,K_p)$ which is totally
ramified.
\end{itemize}
Thus, if $\Lambda$ is an acceptable collection of local specifications with finite total ramification, then $\Lambda_p = \Lambda_p^{\rm ntr}$  for all $p$ large enough. Additionally: 
\begin{enumerate}
\item[\rm(a)]  Let $\cF_4(\Lambda)$ denote the set of pairs $(L, K)$
  consisting of (an isomorphism class of) a quadratic extension $L$ of
  a quadratic field $K$ over $\Q$ such that for every place $v \in \nu(\Q)$, we have
  $$(L_v,K_v):=(L \otimes_{\bQ} \bQ_v, K\otimes_{\bQ} \bQ_v) \in
  \Lambda_v.$$
\item[{\rm (b)}] Let $\FF_2(\Lambda)$ denote the set of quadratic
  extensions $K/\bQ$ such that there exists a quadratic extension $L$
  of $K$ such that $(L,K) \in \cF_4(\Lambda)$.
\item[{\rm (c)}] For a quadratic field $K \in \FF_2(\Lambda)$, let $\FF_4(\Lambda,K)$ be
  the set of $L$ such that $(L,K)\in\FF_4(\Lambda)$.
  \item[{\rm (d)}]   Let $S(\Lambda)\subset\Val(\bQ)$ 
consist of $2$, $\infty$, and all the odd primes $p$ such that
$\Lambda_p \neq \Lambda_p^{\rm ntr}$. (Note that since $\Lambda$ is acceptable with finite total
ramification, the set $S(\Lambda)$ is finite, and that the complement of $S(\Lambda)$ in the set of places of $\Q$ consisting of every odd prime $p$ such that $\Lambda_p = \Lambda_p^{\rm ntr}$. ) 

  \item[{\rm (e)}] For 
$K \in \FF_2(\Lambda)$, let $\cS_{K}(\Lambda)$ denote the valuations of $K$ above the
places in $S(\Lambda)$, and let
$\cS =\cS_K^{\rm ram}(\Lambda)$ denote the union of the set $\cS_K(\Lambda)$ and the set of
ramified places in $K$.
\item[{\rm (f)}] For $K\in\cF_2(\Lambda)$, let $\FF_4^\un(\Lambda,K)
  \subseteq \cF_4(\Lambda,K)$ be the subset of quadratic extensions
  of $K$ that are unramified away from $\cS_K(\Lambda)$.
\end{enumerate}
Finally, we call $(L,K) \in \FF_4(\Lambda)$ a {\em $D_4$-pair} if and only if $L$ is a quartic $D_4$-field. 
\end{defn}

\begin{equation}\label{eqmass}
\Mass^\flat(\Lambda):=\Bigl(\sum_{(L,K)\in\Lambda_\infty}\frac{1}{\#\Aut(L,K)}\Bigr)
\prod_{p\in S(\Lambda)} \Bigl(\sum_{(L,K)\in\Lambda_p}
\frac{1}{\#\Aut(L,K)}\Bigr)\Bigl( 1-\frac{1}{p}\Bigr)^3.
\end{equation}

\begin{thm}\label{thD4quar}
Let $\Lambda$ be an acceptable collection of local
specifications with finite total ramification as in Definition \ref{localspec},   and recall $\Mass^\flat(\Lambda)$ defined in \eqref{eqmass}. Let $N_{D_4}(\Lambda;X)$
denote the number of $D_4$-pairs $(L,K)\in\FF_4(\Lambda)$ such that
$\H(L,K)<X$ (see Definition \ref{sheight}). We then have:
\begin{equation*}
N_{D_4}(\Lambda;X)=\frac{1}{4}\Mass^\flat(\Lambda)\prod_{p\not\in S(\Lambda)}
\Bigl(1+\frac{3}{p}\Bigr)\Bigl(1-\frac{1}{p}\Bigr)^3\cdot X\log^2 X
+o(X\log^2 X).
\end{equation*}
\end{thm}

\begin{thm}\label{thunif}
Let $X>Z>0$ be real numbers. The number of octic $D_4$-fields $M$ such
that $\Delta(M)<X$ and the product of primes at which $M$ totally
ramifies is greater than $Z$ is bounded by $$O_\epsilon\left(\frac{X^{1/4}(\log
X)^2}{Z^{\delta-\epsilon}}\right)$$ for some $\delta > 0$.
\end{thm}

\begin{thm}\label{th:congQuartic}
Let $\Lambda$ be an acceptable collection of quartic local
specifications. Then
\begin{equation*}
N_\Delta(\Lambda;X)=\frac{1}{4}\Bigl(\sum_{(L,K)\in\Lambda_\infty}\frac{1}{\#\Aut(L,K)}\Bigr)
\prod_p\Bigl(\sum_{(L,K)\in\Lambda_p}\frac{|\Delta(M_{(L,K)})|_p^{\frac14}}{\#\Aut(L,K)}\Bigr)
X^{\frac14}\log^2 X^{\frac14}+o(X^{\frac14}\log^2 X),
\end{equation*}
where $M_{(L,K)}$ is the octic algebra corresponding to $(L,K)$.
\end{thm}

\begin{proof}
Since $\Lambda$ is acceptable, it follows that there exists a finite
set $S_\Lambda$ of primes, containing $2$, such that if $p\not\in
S_\Lambda$, then $\Lambda_p$ contains all non-totally ramified pairs.
We assume that $\Lambda$ is a {\it discriminant stable} collection,
i.e., the set $\Lambda_p$ consists of elements whose octic
discriminants have the same valuation $\lambda_p$ for every finite
place $p\in S_\Lambda$. By additivity, the general result clearly
follows from the result on discriminant stable collections, and so
there is no loss of generality in this assumption.  For a positive
integer $T$ not divisible by any primes in $S_\Lambda$, we construct a
collection $\Lambda(T)=(\Lambda(T)_v)_v$ of quartic $D_4$ local
specifications as follows. For $v\in S_\Sigma$ and $p(v) \nmid T$
$\Lambda(T)_v=\Lambda_v$ for all infinite places and for $p \in S_\Sigma$ ; for $p\mid T$, define $\Lambda(T)_p$ to be
the set of all totally ramified extensions; for $p\not\in S_\Sigma$
and $p\nmid T$, define $\Lambda(T)_p$ to be the set of all extensions
which are not totally ramified. Clearly $\Lambda(T)$ is acceptable and
has finite total ramification.

Let $H_T$ denote the $S(\Lambda(T))$-height on $\FF_4(\Lambda(T))$.
Consider a pair $(L,K)\in\FF_4(\Lambda(T))$ and denote the Galois
closure of $L$ over $\bQ$ by $M$. Then we have
\begin{equation*}
|\Delta(M)|=\Bigl(\prod_{p\in S_\Lambda} \lambda_p\Bigr)\cdot T^6\cdot H_T(L,K)^4.
\end{equation*}
Therefore, from Theorem \ref{thD4quar}, we have
\begin{equation*}
\begin{array}{rcl}
\displaystyle N_\Delta(\Lambda(T);X)&=&
\displaystyle N_{D_4}\Bigl(\Lambda;\Bigl(\prod_{p\in
S_\Lambda} \lambda_p^{-1/4} \Bigr)T^{-3/2}X^{1/4}\Bigr)
\\[.2in]&\sim &\displaystyle
\frac14\Mass^\flat(\Lambda(T))\Bigl(\prod_{\substack{p\not\in S_\Lambda\\p\nmid T}}
\Bigl(1+\frac{3}{p}\Bigr)\Bigl(1-\frac{1}{p}\Bigr)^3\Bigr)\cdot
\Bigl(\prod_{p\in S_\Lambda}\lambda_p^{-1/4}\Bigr)T^{-3/2}X^{1/4}\log^2 X^{1/4}.
\end{array}
\end{equation*}
We simplify the above by expanding $\Mass^\flat$ and writing
\begin{equation*}
\Mass^\flat(\Lambda(T))\Bigl(\prod_{p\in S_\Lambda}\lambda_p^{-1/4}\Bigr)T^{-3/2}
=\Bigl(\sum_{(L,K)\in\Lambda_\infty}\frac{1}{\#\Aut(L,K)}\Bigr)
\prod_{p\in S(\Lambda(T))}\Bigl(\sum_{(L,K)\in\Lambda_p}
\frac{|\Delta(M_{(L,K)})|_p^{\frac14}}{\#\Aut(L,K)}\Bigr).
\end{equation*}
The above equality is true since we have
$|\Delta(M_{(L,K)})|_p=\lambda_p^{-1}$ for $p\in S_\Lambda$ and
$(L,K)\in\Lambda_p$, and since for any odd prime $p$, the sum of
$|\Delta(M_{(L,K)})|_p^{\frac14}/\#\Aut(L,K)$ over all totally
ramified pairs is $p^{-3/2}$.

Let $Z>0$ be a real number, and define $\cT(Z)$ to be the set of all
squarefree integers which are the product of primes $p\leq Z$ with
$p\not\in S_\Lambda$. From Theorem \ref{thunif}, we have
\begin{equation*}
N_\Delta(\Lambda;X)=\sum_{T\in\cT(Z)}N_\Delta(\Lambda(T);X)+O_\epsilon
\Bigl(\frac{X^{1/4}\log^2 X}{Z^{\delta-\epsilon}}\Bigr).
\end{equation*}
Since we have
\begin{equation*}
\sum_{(L,K)}\frac{|\Delta(M_{(L,K)})|_p^{\frac14}}{\#\Aut(L,K)}=
\Bigl(1+\frac{3}{p}\Bigr),
\end{equation*}
where the sum is over all non-totally ramified pairs $(L,K)$ over
$\Q_p$, we sum over $T\in\cT(Z)$ and obtain
\begin{equation*}
\begin{array}{rcl}
\displaystyle \frac{N_\Delta(\Lambda;X)}{X^{1/4}\log^2 X^{1/4}}=
\displaystyle\Bigl(\sum_{(L,K)\in\Lambda_\infty}
\frac{1}{\#\Aut(L,K)}\Bigr)
\prod_{p\leq Z}\Bigl(\sum_{(L,K)\in\Lambda_p}
\frac{|\Delta(M_{(L,K)})|_p^{\frac14}}{\#\Aut(L,K)}\Bigr)
\prod_{p>Z}\Bigl(1+\frac{3}{p}\Bigr)\Bigl(1-\frac{1}{p}\Bigr)^3
+O_\epsilon(Z^{-\delta+\epsilon}).
\end{array}
\end{equation*}
Letting $Z$ tend to infinity yields the result.
\end{proof}

\end{document}
